固定 map/reduce 阶段适合规则清楚的大批量流水线。
对账也可能在读到一类差异后，才决定继续展开哪些计算：
取得一个分区图像，拆分其中的付款记录，
对未决记录取得对应证据，再汇合结果。
有些步骤只处理输入并产生输出，
有些步骤要反复使用已经加载的图像或连接。
这要求运行时同时安排短任务与保留状态的执行单元。

Ray 是面向分布式应用的执行框架。
它把远程函数作为任务，把保留实例状态的远程执行单元作为 Actor，
以对象引用连接两者的输入与输出。
程序运行到某个分支时，可以继续提交任务，
计算图因此随执行过程增长。
运行时负责等待依赖、选择节点、传递对象和处理相应重试，
应用负责规定计算及状态变更的含义~\cite{ray-dynamic-design,ray259-objects,ray259-actors}。

从集群看，应用主进程 driver 提交工作，
节点上的 worker 执行函数或驻留 Actor，
raylet 管理本节点的执行资源、worker 与对象传输相关工作。
对象存储承载可共享的中间值，
集群级控制服务保存节点和 Actor 等元数据。
图~\ref{fig:ch20-ray-cluster} 中，调度控制、对象字节与外部业务写入各有路径。
本节按 Ray 2.59.0 解释这些职责。
Ray 统一任务、Actor 和动态依赖的早期设计提供了这一组织思想；
产品演进后的对象 owner 与 GCS 分工，在下面按现行机制分别展开。

\CTwentyFigure{ch20-29}{Ray 的集群控制、节点执行、对象传输与外部权威分别连接}{fig:ch20-ray-cluster}

\subsection{对象引用把依赖交给运行时}

提交远程函数后，调用方先取得 \code{ObjectRef}，
即对未来结果的引用。
它可以继续作为下游任务的参数，
调用方无须先把大对象取回 driver，再重新发送。
对象引用表达哪个结果被依赖，
真正的对象字节则由运行时按执行位置取得。

设任务 A 读取固定图像并产生规范化记录，
任务 B 按付款身份分组，任务 C 汇总差异。
A 的输出引用直接成为 B 的输入，
B 的输出引用成为 C 的输入。
以顶层参数传入的引用，Ray 会在任务执行前取得并解引用为所需值。
若引用藏在列表等嵌套对象中，任务接收的仍是引用，
程序要在需要值时显式取出，或将它继续传给其他任务。
这两种传法分别适合立即计算和只转交引用的步骤~\cite{ray259-objects}。

\CTwentyFigure{ch20-08}{任务提交形成对象依赖，数据无须逐段绕经 driver}{fig:ch20-ray-dependencies}

对象值在 Ray 中按不可变结果使用。
需要改变数据时，任务产生新对象，再把新引用交给后续计算。
A 的固定输出因而可以同时供 $B_1$ 和 $B_2$ 读取，
两者各自产生结果。
共享可变状态则可以放入 Actor 或外部事务域，
由相应实例或事务规则仲裁更新。

对象引用关联运行时中的对象元数据、引用关系和取值路径。
对账需要在 driver 或集群结束后继续保存的输入与输出时，
把对象内容写入可靠存储，
发布第十九章所述带来源、应用位置与摘要的清单，
新的运行时便可从稳定文件身份重新建立引用。

\subsection{对象值、owner 与存储位置}

Ray 对象具有值与元数据两部分。
值可以存在节点对象存储中，位置和引用等元数据由对象 owner 管理。
owner 是创建原始 \code{ObjectRef} 的 worker，
通常是调用远程函数的进程，
与实际计算结果的 worker 可以位于不同机器。
driver 提交 A 时，A 由另一进程计算，
driver 可能同时是 A 输出引用的 owner~\cite{ray259-objects}。

大对象进入节点的共享对象存储。
多个本地 worker 可以共享访问，远端 worker 按需要取得副本。
小对象还可能通过内联路径传递，
所以对象引用与一份独立磁盘文件没有一一对应关系。
对象存储负责值的承载和传输，
owner 负责相关元数据，两者丢失后需要的信息也不同。

\CTwentyFigure{ch20-09}{对象的计算者、owner 与当前值副本可以位于不同进程}{fig:ch20-ray-owner}

内存压力较大时，object spilling 将对象溢写到配置的存储位置，
需要时再读回。
它为对象存储腾出内存，
溢写文件能跨越哪类故障取决于后端及其存活条件。
本机磁盘的溢写与独立存储中的固定作业输入，
分别受节点故障域与外部存储故障域约束。

分布式引用计数跟踪哪些调用方、下游任务或 Actor 持有对象引用，
在引用关系结束后回收相应数据。
应用若不断提交任务并保存所有引用，
driver 元数据、待执行队列和对象存储可能一同增长。
分批提交并通过等待接口取得部分结果，
可以使已经完成且不再引用的对象退出活跃集合。
并行度由 CPU、内存和下游处理能力共同约束。

\subsection{丢失的值怎样重建}

对象值丢失时，Ray 先寻找其他节点上的副本。
没有可用副本时，可以沿 lineage 重执行产生对象的任务，
并递归取得输入。
lineage 是任务与依赖的来源关系，
它说明这些结果由哪个计算、哪些参数产生，
使运行时能够找到重新产生字节的路径~\cite{ray259-objects}。

重建需要输入可取得、对象及其递归依赖由任务产生、
相应 lineage 可用、有关 owner 存活，并且重试预算尚未耗尽。
任务还要满足重建所要求的确定性和副作用条件。
\code{ray.put} 将已有值直接放入运行时，
没有同样的任务生成路径，丢失后无法沿这条计算链重建。
Actor 任务输出默认也不按普通无状态任务自动重建。
将 \code{max\_task\_retries} 配成非零后，
方法可以为重建而再次执行，其状态变化和外部效果要按方法自身规则处理。

\CTwentyFigure{ch20-10}{值副本丢失可以沿任务来源重算，owner 丢失走另一条失败路径}{fig:ch20-ray-reconstruction}

owner 退出时，Ray 当前版本不能从值副本自动恢复原 owner 状态。
后续取值会进入 \code{OwnerDiedError} 等失败路径，
残存值副本也会被清理以回收资源。
需要跨 driver 生命周期恢复的作业，
另存逻辑任务、输入版本和正式发布结果，
新 driver 据此组织计算并创建新的运行时引用。

这也决定业务数据怎样放置。
如果原 $S$ 到 $P/C/T$ 的唯一绑定只在 driver 字典或一个 \code{ObjectRef} 中，
driver 退出就会丢失接续原操作所需的依据。
本书将该绑定保存在复制权威里。
新的计算进程先读取原绑定，再使用新的对象引用承载输入，
计算重启与业务责任恢复由此接上各自的保存来源。

\subsection{调度资格、资源与局部性}

Ray 选择执行位置时，同时考虑输入依赖与资源要求。
节点拥有任务所需资源类型和总量，称为可行。
可行节点的资源若被其他任务占用，当前便不可用。
只有 CPU 的节点对于要求 GPU 的任务不可行，
等待 CPU 空闲不会改变其资源类型。
增加可行节点或调整任务要求后，调度者才有新的选择~\cite{ray259-scheduling}。

默认调度结合逻辑资源利用率与分布策略选择可行节点。
普通任务的大参数局部性还会影响选择，
优先考虑已经持有较多输入字节的可用节点。
指定 \code{SPREAD} 时，调度者尽量分散任务或 Actor。
placement group 把若干资源束按指定关系一起放置，
用于组织需要配合运行的一组执行单元。
节点亲和表达对特定节点的偏好或要求。
这些策略具有各自的优先关系，
普通任务的参数局部性规则也没有普遍扩展到 Actor。

逻辑 CPU 或 GPU 数量是调度配额。
进程申请一个 CPU 后，还可能启动多个操作系统线程，
应用和底层数值库要另外设置合适的线程数。
声明零 CPU 的任务也会实际消耗计算资源。
将声明配额与实际线程、内存使用对应起来，
调度者保存的资源视图才接近机器的实际承载。

\CTwentyFigure{ch20-11}{依赖就绪、资源可行与执行位置共同决定任务何时运行}{fig:ch20-ray-placement}

GCS，即 Global Control Service，管理集群级元数据，
包括节点、Actor 与 placement group 等信息。
对象 owner 保存某个对象的关键元数据，
GCS 保存集群层面的控制信息，
应用权威保存业务绑定和原结果，三者有不同的恢复来源。
默认 GCS 将数据保存在内存中，其失败会使集群失败。
要使它恢复，需要配置可重新加载的持久后端，
并安排故障检测、重启及单一有效 GCS~\cite{ray259-gcs}。

Ray 2.59.0 文档给出外部 Redis 与嵌入式 RocksDB 两种后端。
外部 Redis 的正式支持范围限定在文档列出的 KubeRay、Ray Serve 配置中。
嵌入式 RocksDB 后端处于 alpha 阶段，只支持 Linux，
数据所在持久卷要能在新 head 上重新挂载。
恢复期间，Actor 创建与重建、资源管理和 worker 注册等功能可能暂停。
已经运行的任务和 Actor、已有对象按各自存活条件继续，
新的调度能力则等待控制服务恢复。

\subsection{任务重试的默认含义}

普通任务的 worker 意外退出时，Ray 默认最多重试 3 次，
\code{max\_retries} 可以调整预算。
应用代码抛出异常时，默认不自动重试，
异常作为 \code{RayTaskError} 传给读取结果或依赖该结果的调用方。
配置 \code{retry\_exceptions} 后，
运行时才按所选异常类型与重试预算重执行~\cite{ray259-task-retry}。

一次任务可能完成外部写入后，在返回结果时退出。
无论运行时将它识别为 worker 失败，
还是应用将错误列为可重试异常，
重执行都可能再次到达同一个外部接口。
读取固定对象并产生确定结果的纯计算，重执行主要增加计算成本。
付款调用则从输入与权威取得原 \code{scope/pay-C}，
平台识别原操作，权威保存原 $C/T$ 的完整结果。

预算耗尽表示运行时停止为这次任务继续尝试。
已经形成的未决付款由业务恢复者从原 $C/T$ 接续查询。
业务已保存原结果时，下游计算失败只改变该计算的完成状态，
权威中已经成功的购买结果继续供原身份读取。
运维据此分别观察任务状态、业务结果与当前恢复进度。
